%!TEX encoding = UTF-8 Unicode

%----------------------------------------------------------------------------------------
%   CHAPTER 5
%   Translator: laserdog
%   Proofreder: SI, cancan123456
%----------------------------------------------------------------------------------------

\chapterimage{chapter_head_1.pdf} % Chapter heading image

\chapter[测量]{Measuring Nature 测量}\label{chap5}

现在我们已经发现对称性与守恒量之间的关系了，接下来就可以利用这种联系。用更专业的话说，Noether定理建立了对称变换的生成元与相应守恒量之间的联系。本章将利用这种联系。

守恒量常常被物理学家们用来描述物理系统，因为无论这个系统经历了怎样复杂的变化，守恒量都是一样的。比如，物理学家们会使用动量，能量或角动量描述一个实验。Noether定理提示了一个至关重要的思想：

描述自然的物理量等同于其对应的生成元：

\begin{align}\label{equ5.1}
\text{物理量}\Rightarrow\text{相应的对称性的生成元}
\end{align}

我们会看到这种对应会自然地导出量子理论。下面我们考虑粒子理论，使这句话更为具体。

\section[量子力学中的算符]{Operators of Quantum Mechanics 量子力学中的算符}\label{sec5.1}

拉格朗日量在空间平移生成元的作用下不变，这种不变性可导出动量守恒。由此可得\footnote{译者注：这里再次提醒一下，本书中虚数单位与指标经常共用字母$i$，相信聪明的读者有能力分辨。}
\[\text{动量} \hat{p}_i\to\text{空间平移操作的生成元：} - i\partial_i \]

{\itshape 惯例用一个尖帽子$\hat{O}$来表示一个算符。}

类似的，时间平移生成元带来的不变性给出能量守恒，也就是
\[\text{能量} \hat{E}\to\text{时间平移操作的生成元} i\partial_0 \]

没有对应于“位置守恒”的对称性，所以位置并没有伴随生成元，我们有\mpar{或者说，我们可以观察对应于推动不变性的守恒量。注意，我们在 \ref{sec4.6} 节得到的守恒量对应于非相对论性 Galileo 推动，因为我们的推导是从非相对论性的拉格朗日量出发的。尽管如此，我们可以对相对论性拉格朗日量进行相同的推导，得到守恒量$t\vec{p}-\vec{x}E$。相对论性能量由$E=\sqrt{m^2+p^2}$给出。在非相对论性极限$c\to\infty$下，我们有$E\approx m$，并且 Lorentz 推动下的守恒量退化为 Galileo 推动的守恒量。因此，粒子理论的守恒量为$M_i = (tp_i-x_iE)$。推动的生成元（见\eqref{equ3.248}，其中$K_i = M_{0i}$）为$K_i = i(x_0\partial_i-x_i\partial_0)$。比较二式，其中$x_0=t$，可得$M_i = (tp_i-x_iE)\leftrightarrow K_i = i(t\partial_i-x_i\partial_0)$。因此，利用前文提及的对应关系，我们可以直接得到：位置 $\hat{x}_i\to x_i$.}
\[\text{位置} \hat{x}_i\to x_i \]

上面利用（微分）算符给出了描述自然的理论中使用的物理量。接下来我们自然要问：这些算符作用在什么上面？这些算符与实验中的可观测量有什么联系？下一章会仔细讨论这件事情，在这里只需要知道物理量，也就是算符，是作用在{\itshape 某物}上的。不妨就认为它作用于一个抽象的对象上，称之为$\Psi$，后面会探究这个“{\itshape 某物}”是什么鬼。

现在可以导出一些非常重要的\mpar{如果你完全不了解量子力学，你也许会感到奇怪：为什么这些简单的式子会如此重要？但也许你听说过Heisenberg不确定性原理。在第\ref{sec8.3}节，我们会更进一步地研究量子力学的形式与结构，然后我们就会发现：这个方程表明我们不能以任意精度同时测量一个粒子的动量和位置。我们将物理量解释为测量算符，而这个方程告诉我们先测量动量再测量位置与先测量位置再测量动量是完全不一样的。}量子力学的方程了。就像上面已经解释过的那样，我们假定算符作用在抽象的$\Psi$上。于是有
\begin{align}\label{equ5.2}
\begin{split}
[\hat{p}_i, \hat{x}_j] \Psi &= (\hat{p}_i \hat{x}_j - \hat{x}_j \hat{p}_i) \Psi = i(-\partial_i \hat{x}_j + \hat{x}_j \partial_i) \Psi \\
&\underbrace{=}_{\mathclap{\text{莱布尼兹律}}} -i(\partial_i \hat{x}_j) \Psi - \cancel{i\hat{x}_j (\partial_i \Psi)} + \cancel{i\hat{x}_j (\partial_i \Psi)} \underbrace{=}_{\mathclap{\partial_i \hat{x}_j = \frac{\partial x_j}{\partial x_i}}} -i \delta_{ij} \Psi
\end{split}
\end{align}

这个方程对任意的$\Psi$都成立，因为我们没有对$\Psi$做出任何假设，因此我们可以不考虑它\mpar{如果你以前没有过这个观念，注意我们可以把{\bfseries 每}一个矢量方程都写成只包含数的形式。比如，牛顿第二定律$\vec{F} = m\ddot{\vec{x}}$ 可以写成 $\vec{F}\vec{C} = m\ddot{\vec{x}}\vec{C}$，显然对于任意矢量 $\vec{C}$ 该等式都成立。然而，如果该等式对于任意 $\vec{C}$ 都成立，一直带着 $\vec{C}$ 没什么意义。}，从而直接写出
\begin{align}\label{equ5.3}
[\hat{p}_i,\hat{x}_j] = -i\delta_{ij}.
\end{align}

\subsection[自旋与角动量]{Spin and Angular Momentum 自旋与角动量}\label{sec5.1.1}

上一章的第\ref{sec4.5.4}节中讲过，旋转不变性带来的守恒量包含两部分，第二部分是轨道角动量，对应于{\bfseries 无限}维生成元表示。
\[\text{轨道角动量}\hat{L}_i\to\text{转动生成元（无限维表示）}i\tfrac{1}{2}\epsilon_{ijk}(x^j\partial^k-x^k\partial^j)\]
类似的， 第一部分称作自旋，对应{\bfseries 有限}维度生成元\mpar{这一部分是场分量混合的不变性带来的守恒量，因此是有限维的表示。}的表示。
\[\text{自旋}\hat{S}_i\to\text{转动生成元（有限维表示）}S_i \]
如前面\eqref{equ4.30}式解释过的那样，$S_{\mu\nu}$和转动算符生成元$S_i$的关系为$S_i = \tfrac{1}{2}\epsilon_{ijk}S_{jk}$。

比如，当我们考虑自旋$\tfrac{1}{2}$场的时候，我们需要用到我们在\ref{sec3.7.5}节中得到的二维表示：
\begin{align}\label{equ5.4}
\hat{S}_i = i\frac{\sigma_i}{2}
\end{align}
其中$\sigma_i$是Pauli矩阵。在学习如何利用本章导出的这些算符之后，我们将在\ref{sec8.5.5}节深入讨论这种有趣而奇怪的角动量。时刻注意，只有这两部分角动量的{\bfseries 和}是守恒的。

\section[量子场论的算符]{The Operators of Quantum Field Theory 量子场论的算符}\label{sec5.2}

场论的最重要的客体当然就是场。场是{\bfseries 时空位置的函数}\mpar{这里，$x = x_0,x_1,x_2,x_3$，包含时间$x_0=t$。}$\Phi = \Phi(x)$。接下来，我们希望描述时空中各点的相互作用，因此我们需要考虑动力学变量$\pi = \pi(x)$的密度，而不是其总量$\Pi = \int d^3x\pi(x)\neq\Pi(x)$。

我们上一章发现了场自身的平移不变性$\Phi\to\Phi-i\epsilon$产生了一种新的守恒量，称为共轭动量$\Pi$。类似于上一节给出的对应关系，我们用对应的生成元\eqref{equ4.60}式来表示共轭动量密度：
\[\text{共轭动量密度}\pi(x)\to\text{场自身平移的生成元：}-i\frac{\partial}{\partial\Phi(x)}.\]
类似粒子理论，我们可以继续找出动量、角动量、能量及其对应的生成元。但是，后面会看到量子场论与粒子理论稍稍不同，因此现在的处理已经足够了。

就像上一节讨论的那样，我们需要给算符一些作用的对象。这里同样使用抽象的$\Psi$，后面的章节会具体说明。同样可以得到一个针对量子场论\mpar{最后一步推导中，我们使用了 $\tfrac{\partial x_i}{\partial x_j} = \delta_{ij}$ 的类比，以给出delta分布$\tfrac{\partial f(x_i)}{\partial f(x_j)} = \delta(x_i-x_j)$，我们可以更严格地证明该等式。更具体的内容请看附录D.2。}十分重要的方程。
\begin{align}\label{equ5.5}
\begin{gathered}
[\Phi(x),\pi(y)]\Psi = \left[\Phi(x),-i\frac{\partial}{\partial\Phi(y)}
\right]\Psi\\
\underbrace{=}_{\text{莱布尼兹律}}\cancel{-i\Phi(x)\frac{\partial\Psi}{\partial\Phi(y)}} + \cancel{i\Phi(x)\frac{\partial\Psi}{\partial\Phi(y)}} + i\left(\frac{\partial\Phi(x)}{\partial\Phi(y)}\right)\Psi = i\delta(x-y)\Psi
\end{gathered}
\end{align}

同样的，由于这个式子对任意的$\Psi$都成立，可以写成：
\begin{align}\label{equ5.6}
[\Phi(x),\pi(y)] = i\delta(x-y).
\end{align}
类似的，对于场有多于一个分量的情况，我们有
\begin{align}\label{equ5.7}
[\Phi_i(x),\pi_j(y)] = i\delta(x-y)\delta_{ij}.
\end{align}

后面我们会看到，量子场论的几乎所有内容都要从这个简单的方程出发。
